% ==================== 第7章 消融实验 ====================
\section{消融实验}

本章的消融在锚点分支和最终系统上分别进行。锚点分支的损失只用于比较证据策略，不与第\ref{ch:empirical}章的融合主表直接比较；最终系统的消融使用与主表相同的后验和读出。

\subsection{锚点提示词与时间演化}

文本锚点的作用首先出现在纯LLM分支。旧提示词使用结构化摘要，新提示词使用时点原文、段落编号和逐字引文。DeepSeek上的准确率、平衡准确率和宏F1分别从.445、.307、.242提高到.502、.397、.364。来源落地和逐字引文让模型能够区分“指引区间”“历史实际”和“分析师预测”，减少了把累计增速当作全年增速的错误。

时间演化的作用较小但方向一致。新锚点从静态可靠度切换到按揭示日更新的Beta可靠度，损失从2.908降到2.855；旧锚点从2.899降到2.793。准确率变化不大，演化主要改善概率校准和决策损失，而不是点分类。这个结果符合演化的设计目标：它更新的是锚点类型的可靠性，不是标签本身。

\subsection{资格判定}

表\ref{tab:qual-ablation}比较完整资格判定与接受全部解析锚点。完整判定拒绝空泛、过期、对抗和不合格的锚点，接受率在三个基座上分别为27.6\%、42.7\%和68.3\%；接受全部条件保留所有解析结果。

\begin{table}[htbp]
\centering
\zihao{5}
\setlength{\tabcolsep}{4pt}
\caption{资格判定消融（227条核验断言，锚点分支）。差值为接受全部减完整判定，正值表示完整判定损失更低。}
\label{tab:qual-ablation}
\begin{tabular}{lrrrr}
\toprule
基座 & 完整判定 & 接受全部 & 损失差 [95\% CI] & 接受率 \\
\midrule
DeepSeek & 3.304 & 3.507 & $+0.203$ $[-0.106,0.520]$ & 27.6\% \\
GLM & 3.339 & 3.427 & $+0.088$ $[-0.176,0.361]$ & 42.7\% \\
Qwen3 & 4.599 & 4.652 & $+0.053$ $[-0.141,0.256]$ & 68.3\% \\
\bottomrule
\end{tabular}
\end{table}

三个基座的点估计都支持完整判定，DeepSeek的增益最大，Qwen3的增益最小。三个配对区间都包含零，因此资格判定的作用按方向性证据报告。Qwen3的接受率明显高于另外两个基座，说明其抽取器产生了更多边缘锚点，资格判定拒绝的比例更高，但剩余锚点的区分力有限。

\subsection{融合权重}

表\ref{tab:weight}给出验证集上不同融合权重的决策损失。权重$w$是锚点后验在融合后验中的占比，$w=0$对应纯随机森林，$w=1$对应纯锚点后验。DeepSeek在$w=0.2$取得最低验证损失，GLM和Qwen3在$w=0.1$取得最低验证损失；测试集使用验证集选择的权重，不参与调参。

\begin{table}[htbp]
\centering
\zihao{5}
\setlength{\tabcolsep}{4pt}
\caption{融合权重的验证集损失（185条验证断言）。粗体为验证集选择。}
\label{tab:weight}
\begin{tabular}{lrrrrrrr}
\toprule
基座 & 0.0 & 0.05 & 0.1 & 0.2 & 0.3 & 0.5 & 1.0 \\
\midrule
DeepSeek & 3.341 & 3.308 & 3.265 & \textbf{3.232} & 3.341 & 3.503 & 3.859 \\
GLM & 3.341 & 3.211 & \textbf{3.200} & 3.286 & 3.362 & 3.416 & 3.708 \\
Qwen3 & 3.341 & 3.362 & \textbf{3.330} & 3.686 & 3.859 & 4.162 & 6.519 \\
\bottomrule
\end{tabular}
\end{table}

三个基座的最优权重都偏向随机森林，DeepSeek为0.2，GLM和Qwen3为0.1。锚点后验单独使用时损失明显更高，Qwen3在$w=1$时达到6.519。融合的价值来自补充信号，而不是替代结构化后验。

\subsection{类型库的验证集扩展}

类型库扩展是算法的一个训练步骤，不是独立的预测模块。算法\ref{alg:typeexpand}先在验证集上挖掘高损失断言，由LLM提出候选类型，再在选择集上重抽锚点并计算损失；只有验证损失下降超过$\delta$的候选才进入类型库。接受后的类型库冻结，测试集只评估一次。


DeepSeek上的候选类型库包含“指引中点与当前季度收入之比”和“财年周数效应”两个类型。验证选择集上，基础类型库的损失为3.943，候选类型库为3.792，下降0.151，超过阈值，候选被接受。提示词单独修改和类型加提示词组合两个候选没有通过验证门控。冻结类型库后，测试集上的准确率从.555提高到.577，宏F1从.513提高到.530，损失从2.934降到2.907；配对损失差为$-0.026$ $[-0.159,0.097]$，区间包含零。

表\ref{tab:typeevolve}汇总这一过程。类型库扩展的测试增益是小的，但方向与验证集一致，且没有使用测试标签。GLM和Qwen3没有候选通过验证门控，保留基础类型库，因此主表中只有DeepSeek使用扩展类型库。

\begin{table}[htbp]
\centering
\zihao{5}
\setlength{\tabcolsep}{4pt}
\caption{DeepSeek类型库扩展（227条核验断言）。Loss为全动作决策损失。}
\label{tab:typeevolve}
\begin{tabular}{lrrrr}
\toprule
系统 & Acc. & Bal. & Macro & Loss \\
\midrule
基础类型库 & .555 & .511 & .513 & 2.934 \\
扩展类型库 & .577 & .526 & .530 & 2.907 \\
差值 & $+.022$ & $+.015$ & $+.017$ & $-0.026$ $[-0.159,0.097]$ \\
\bottomrule
\end{tabular}
\end{table}

\subsection{去掉锚点结构的对照}
\label{sec:no-anchor}

前面的消融检验的是组件内部，例如资格判定是否比接受全部更好。一个更根本的问题是：如果去掉三层锚点结构，只保留时点证据、LLM标签分数、温度缩放和全动作读出，损失会差多少。这个对照直接检验锚点结构本身的必要性。

本文从同一批锚点抽取响应中只保留标签分数，去掉全部锚点对象、资格判定和Beta可靠性，在验证集上拟合温度参数，再在测试集上使用同一个全动作读出。这个对照可以看作FinAbstain式选择性预测在本任务上的近似：时点证据、概率输出、温度缩放和损失感知弃权都在，但没有多类别资格判定和类型可靠性演化。表\ref{tab:noanchor}给出结果。去掉锚点结构后，DeepSeek、GLM和Qwen的损失分别上升0.379、0.317和0.899，准确率分别下降5.3、8.8和13.7个百分点。时点证据和选择性输出仍然存在，性能却明显低于完整系统，说明增益不能只归因于“时点检索加弃权”这一一般性组合。

\begin{table}[htbp]
\centering
\zihao{5}
\setlength{\tabcolsep}{4pt}
\caption{去掉三层锚点结构的架构级对照（227条核验断言）。温度在验证集上选择，测试集只评估一次。}
\label{tab:noanchor}
\begin{tabular}{lrrrrrrrrr}
\toprule
\multirow{2}{*}{方法} & \multicolumn{3}{c}{DeepSeek} & \multicolumn{3}{c}{GLM} & \multicolumn{3}{c}{Qwen3} \\
\cmidrule(lr){2-4}\cmidrule(lr){5-7}\cmidrule(lr){8-10}
 & Acc. & Macro & Loss & Acc. & Macro & Loss & Acc. & Macro & Loss \\
\midrule
无锚点结构 & .524 & .398 & 3.286 & .463 & .347 & 3.260 & .405 & .257 & 3.824 \\
\textbf{FinPlanRAG} & \textbf{.577} & \textbf{.530} & \textbf{2.907} & \textbf{.551} & \textbf{.504} & \textbf{2.943} & \textbf{.542} & \textbf{.501} & \textbf{2.925} \\
\bottomrule
\end{tabular}
\end{table}

这个对照也解释了为什么本文把贡献放在证据资格判定和类型可靠性演化上，而不是放在“选择性预测”上。选择性预测本身是一个一般性工具，FinAbstain等工作已经把它用于金融预测；本文的增量在于锚点结构把证据拆成可判定、可追溯、可按类型更新的对象。去掉这层结构，系统仍然可以选择弃权或裁决集，但失去了区分证据失效模式的能力。

\subsection{数值桥}

数值桥的作用在A股回测上最清楚。DeepSeek完整系统的损失为2.179，去桥变体为2.421，差值为$-0.242$；准确率从.638提高到.783，宏F1从.590提高到.776。桥接把“增速可能高于区间”的标签判断转化为“按当前季度年化，全年增速约为多少”的数值计算，减少了纯标签后验的过度平滑。

桥接的收益依赖抽取质量。GLM完整系统与去桥变体完全相同，说明该基座没有产生足够的数值锚点，桥接未触发；Qwen3的结果见第\ref{ch:empirical}章。数值桥不是无条件增益，它要求抽取器输出可核验的数值字段、单位和期间。对抽取质量较弱的基座，标签后验仍是更稳的回退。

\subsection{消融结论}

消融结果可以归纳为四点。第一，去掉三层锚点结构后，三个基座的损失分别上升0.379、0.317和0.899，说明锚点结构本身提供了不可替代的增量，而不是只有选择性输出在起作用。第二，输出形式和资格判定是稳定的组件：全动作读出在三个基座上显著降低损失，资格判定方向一致但区间包含零。第三，类型库扩展和数值桥是条件增益：它们在DeepSeek上有效，在抽取质量较弱的基座上不触发或收益有限。第四，融合权重偏向随机森林，锚点层提供的是补充信号，而不是替代结构化后验。这个结论与主表的定位一致：框架的优势来自证据状态、决策接口和可审计性，而不是单一组件的普遍提升。
